\chapter{Stress Tensor Trace and Conformal Currents}
\label{ch:volume-ii-stress-tensor-trace-conformal-currents}

\section{The Stress Tensor from Local Translations}

Work in flat \(D\)-dimensional Euclidean spacetime, with coordinates \(x^\mu\)
and metric \(\delta_{\mu\nu}\).  A constant translation by a vector
\(\delta a^\mu\) acts on a scalar field by
\[
  \phi'(x')=\phi(x),
  \qquad
  x'^\mu=x^\mu+\delta a^\mu .
\]
The corresponding active variation at the same coordinate \(x\) is
\[
  \delta\phi(x)=\phi'(x)-\phi(x)
  =
  -\delta a^\mu\,\partial_\mu\phi(x).
\]
The stress tensor is obtained by promoting \(\delta a^\mu\) to a smooth
compactly supported function \(\delta a^\mu(x)\).  For a multiplet of fields
one is free to add to the local variation a term proportional to
\(\partial_\mu\delta a_\nu\):
\[
  \delta\phi(x)
  =
  -\delta a^\mu(x)\,\partial_\mu\phi(x)
  +
  \partial_\mu\delta a_\nu(x)\,\mathcal D^{\mu\nu}\phi(x).
\]
Here \(\mathcal D^{\mu\nu}\) is a fixed linear differential operator on the
fields.  Different choices of \(\mathcal D^{\mu\nu}\) change the stress tensor
by equation-of-motion terms and improvement terms.  Once the convention is
chosen, define \(T^{\mu\nu}\) by
\[
  S[\phi+\delta\phi]
  =
  S[\phi]
  -
  \int\dd^D x\,
  \partial_\mu\delta a_\nu(x)\,T^{\mu\nu}(x),
\]
for compactly supported \(\delta a^\mu(x)\), up to terms proportional to the
equations of motion.  If \(\delta a^\mu\) is constant, the derivative vanishes,
which is the Noether statement of translation invariance.

\begin{figure}[H]
  \centering
  \resizebox{0.94\textwidth}{!}{%
  \begin{tikzpicture}[x=1cm,y=1cm,every node/.style={font=\scriptsize}]
    \node[font=\small] at (-3.4,1.8) {constant translation};
    \draw[-{Latex[length=2mm]}, thick] (-5.1,0.65)--(-3.8,0.65);
    \node[above] at (-4.45,0.65) {\(\delta a^\mu\)};
    \node[align=center] at (-4.45,-0.1)
      {\(\partial_\mu\delta a_\nu=0\)\\\(\delta S=0\)};

    \node[font=\small] at (1.7,1.8) {localized translation};
    \draw[thick, rounded corners] (0.0,0.05) rectangle (3.55,1.15);
    \draw[-{Latex[length=2mm]}, thick] (0.35,0.55)
      .. controls (1.0,1.05) and (1.9,1.05) .. (2.75,0.65);
    \node[above] at (1.7,1.07) {\(\delta a^\mu(x)\)};
    \node[align=center] at (1.75,-0.45)
      {\(\displaystyle
      \delta S=-\int\dd^D x\,
      \partial_\mu\delta a_\nu\,T^{\mu\nu}\)};

    \draw[-{Latex[length=2mm]}, thick] (-1.1,0.55)--(-0.1,0.55);
    \node[align=center, blue!70!black] at (-0.55,1.18)
      {localize};
  \end{tikzpicture}
  }
  \caption{The stress tensor is the coefficient of the derivative of a local
  translation parameter.  Constant translations are recovered by setting
  \(\partial_\mu\delta a_\nu=0\).}
  \label{fig:volume-ii-local-translation-stress-tensor}
\end{figure}

\section{Dilatations and Engineering Dimensions}

A constant infinitesimal scale transformation is the special local
translation
\[
  \delta a^\mu(x)=\sigma x^\mu,
\]
where \(\sigma\) is a constant.  The stress-tensor definition gives
\[
  \delta S
  =
  -\sigma\int\dd^D x\,T^\mu{}_\mu(x).
\]
Thus the trace is the local operator that measures the response to an
infinitesimal dilatation in the chosen stress-tensor convention.

For a free massless scalar,
\[
  S_{\rm kin}[\phi]
  =
  \frac12\int\dd^D x\,\partial_\mu\phi\,\partial^\mu\phi,
\]
scale invariance of the kinetic term fixes the engineering field dimension to
\[
  \Delta_\phi^{\rm eng}=\frac{D-2}{2}.
\]
Equivalently, with \(x'^\mu=(1+\sigma)x^\mu\),
\[
  \phi'(x')
  =
  (1+\sigma)^{-(D-2)/2}\phi(x),
\]
or, at fixed \(x\),
\[
  \delta\phi(x)
  =
  -\sigma x^\mu\partial_\mu\phi(x)
  -
  \sigma\frac{D-2}{2}\,\phi(x).
\]

Let now \(D=d-\epsilon\), and let \(O_I\) be a local monomial or polynomial in
fields and derivatives with engineering dimension
\[
  d_I(\epsilon)=d-\epsilon+\delta_I(\epsilon)=D+\delta_I(\epsilon).
\]
Under the same scale transformation,
\[
  O_I'(x')
  =
  (1+\sigma)^{-d_I(\epsilon)}O_I(x).
\]
For an action of the form
\[
  S[\phi]
  =
  S_{\rm kin}[\phi]
  +
  \int\dd^D x\,\sum_I g_I^\epsilon\,O_I(x),
\]
the interaction part transforms as
\[
  \int\dd^D x'\sum_I g_I^\epsilon O_I'(x')
  =
  \int\dd^D x\sum_I
  g_I^\epsilon(1+\sigma)^{-\delta_I(\epsilon)}O_I(x).
\]
To first order in \(\sigma\),
\[
  \delta S
  =
  -\sigma\int\dd^D x\,
  \sum_I\delta_I(\epsilon)g_I^\epsilon O_I(x).
\]
Comparing this with the stress-tensor definition gives the bare trace identity
\[
  T^\mu{}_\mu
  =
  \sum_I\delta_I(\epsilon)\,g_I^\epsilon\,O_I,
\]
modulo total derivatives and equation-of-motion operators.

\begin{figure}[H]
  \centering
  \resizebox{0.92\textwidth}{!}{%
  \begin{tikzpicture}[x=1cm,y=1cm,every node/.style={font=\scriptsize}]
    \node[font=\small] at (-3.2,1.75) {dilatation};
    \node[draw=black, rounded corners, align=center, inner sep=6pt]
      at (-3.2,0.55)
      {\(\delta a^\mu=\sigma x^\mu\)\\
       \(\delta S=-\sigma\int T^\mu{}_\mu\)};

    \node[font=\small] at (1.9,1.75) {operator dimensions};
    \node[draw=black, rounded corners, align=center, inner sep=6pt]
      at (1.9,0.55)
      {\(O_I'(x')=(1+\sigma)^{-D-\delta_I}O_I(x)\)\\
       \(\delta S=-\sigma\int\sum_I\delta_I g_I^\epsilon O_I\)};

    \draw[-{Latex[length=2mm]}, thick] (-1.2,0.55)--(-0.15,0.55);
    \node[blue!70!black, align=center] at (-0.65,-0.28)
      {compare\\local densities};
    \node[font=\small] at (-0.65,-1.0)
      {\(\displaystyle T^\mu{}_\mu=\sum_I\delta_I(\epsilon)g_I^\epsilon O_I\)};
  \end{tikzpicture}
  }
  \caption{The bare trace is the engineering-dimension defect of the regulated
  interaction terms.}
  \label{fig:volume-ii-scale-trace-engineering}
\end{figure}

\section{The Renormalized Trace Identity}

Minimal subtraction expresses each bare coupling as
\[
  g_I^\epsilon
  =
  \mu^{-\delta_I(\epsilon)}
  \left[
    \lambda_I(\mu)
    +
    \sum_{n\ge1}
    \epsilon^{-n}K_I^{(n)}(\lambda(\mu))
  \right],
\]
where
\[
  \delta_I(\epsilon)
  =
  \delta_I^{(0)}+\epsilon\delta_I^{(1)}.
\]
Bare couplings are independent of \(\mu\).  If
\[
  \beta_I^\epsilon(\lambda)
  =
  \frac{\dd\lambda_I}{\dd\log\mu}
  =
  \beta_I(\lambda)+\epsilon\delta_I^{(1)}\lambda_I
\]
is the \(D=d-\epsilon\) beta function in MS coordinates, then
\[
  0
  =
  \frac{\dd g_I^\epsilon}{\dd\log\mu}
  =
  -\delta_I(\epsilon)g_I^\epsilon
  +
  \sum_J
  \beta_J^\epsilon
  \frac{\partial g_I^\epsilon}{\partial\lambda_J}.
\]
Hence
\[
  \delta_I(\epsilon)g_I^\epsilon
  =
  \sum_J
  \bigl(\beta_J+\epsilon\delta_J^{(1)}\lambda_J\bigr)
  \frac{\partial g_I^\epsilon}{\partial\lambda_J}.
\]
Substituting this relation into the bare trace identity gives
\[
  T^\mu{}_\mu
  =
  \sum_J
  \bigl(\beta_J+\epsilon\delta_J^{(1)}\lambda_J\bigr)
  \frac{\partial}{\partial\lambda_J}
  \left(\sum_I g_I^\epsilon O_I\right).
\]
The derivative in parentheses is the renormalized operator associated with
the MS coordinate \(\lambda_J\):
\[
  [O_J]_\mu
  =
  \frac{\partial}{\partial\lambda_J}
  \left(\sum_I g_I^\epsilon O_I\right).
\]
Taking the finite \(\epsilon\to0\) limit gives the renormalized trace identity
\[
  T^\mu{}_\mu
  =
  \sum_J\beta_J(\lambda)\,[O_J]_\mu .
\]
This identity is an operator identity in flat space for separated insertions;
contact terms, total derivatives, and equation-of-motion operators are fixed
by the corresponding Ward identities and by the chosen stress-tensor
improvement convention.

\begin{figure}[H]
  \centering
  \resizebox{0.95\textwidth}{!}{%
  \begin{tikzpicture}[x=1cm,y=1cm,every node/.style={font=\scriptsize}]
    \node[draw=black, rounded corners, align=center, inner sep=6pt]
      at (-4.2,0.7)
      {\(\displaystyle T^\mu{}_\mu
      =\sum_I\delta_I(\epsilon)g_I^\epsilon O_I\)\\
       bare engineering trace};
    \node[draw=black, rounded corners, align=center, inner sep=6pt]
      at (0,0.7)
      {\(\displaystyle
      \delta_I g_I^\epsilon
      =
      \sum_J\beta_J^\epsilon
      \partial_{\lambda_J}g_I^\epsilon\)\\
       \(\dd g_I^\epsilon/\dd\log\mu=0\)};
    \node[draw=black, rounded corners, align=center, inner sep=6pt]
      at (4.3,0.7)
      {\(\displaystyle
      T^\mu{}_\mu=\sum_J\beta_J[O_J]_\mu\)\\
       finite renormalized trace};
    \draw[-{Latex[length=2mm]}, thick] (-2.35,0.7)--(-1.55,0.7);
    \draw[-{Latex[length=2mm]}, thick] (1.7,0.7)--(2.55,0.7);
    \node[blue!70!black] at (0,-0.8)
      {\([O_J]_\mu=\partial_{\lambda_J}(\sum_Ig_I^\epsilon O_I)\)};
  \end{tikzpicture}
  }
  \caption{The MS trace identity is obtained by converting the engineering
  dimension factor of a bare coupling into the beta-function vector field
  acting on the same bare interaction.}
  \label{fig:volume-ii-renormalized-trace-identity}
\end{figure}

\section{Fixed Points and Conformal Currents}

An RG fixed point in the MS coordinate system is a point \(\lambda_\ast\) at
which
\[
  \beta_J(\lambda_\ast)=0
  \qquad\text{for all }J .
\]
At such a point the beta-function contribution to the trace vanishes:
\[
  T^\mu{}_\mu=0,
\]
again modulo total derivatives, contact terms, and equation-of-motion
operators.  Suppose that the stress tensor has been chosen symmetric,
conserved, and traceless as a flat-space renormalized operator:
\[
  T^{\mu\nu}=T^{\nu\mu},
  \qquad
  \partial_\mu T^{\mu\nu}=0,
  \qquad
  T^\mu{}_\mu=0.
\]
Let \(\xi^\mu(x)\) be a conformal Killing vector field in \(d\) dimensions,
meaning
\[
  \partial_\mu\xi_\nu+\partial_\nu\xi_\mu
  -
  \frac{2}{d}\delta_{\mu\nu}\,\partial_\rho\xi^\rho
  =
  0.
\]
Define the current
\[
  j^\mu_\xi(x)=T^{\mu\nu}(x)\xi_\nu(x).
\]
Its divergence is
\[
  \partial_\mu j^\mu_\xi
  =
  T^{\mu\nu}\partial_\mu\xi_\nu
  =
  \frac12 T^{\mu\nu}
  \bigl(\partial_\mu\xi_\nu+\partial_\nu\xi_\mu\bigr)
  =
  \frac{1}{d}T^\mu{}_\mu\,\partial_\rho\xi^\rho
  =
  0.
\]
Thus each conformal Killing vector gives a conserved Noether current.

The elementary conformal Killing vector fields are
\[
  \xi^\mu(x)
  =
  \begin{cases}
    a^\mu, & \text{translations},\\
    \omega^\mu{}_\nu x^\nu,\quad \omega_{\mu\nu}=-\omega_{\nu\mu},
      & \text{rotations},\\
    c\,x^\mu, & \text{dilatations},\\
    b_\rho(2x^\rho x^\mu-x^2\delta^{\rho\mu}),
      & \text{special conformal transformations}.
  \end{cases}
\]

\begin{figure}[H]
  \centering
  \resizebox{0.96\textwidth}{!}{%
  \begin{tikzpicture}[x=1cm,y=1cm,every node/.style={font=\scriptsize}]
    \node[font=\small] at (-3.8,2.0) {stress tensor data};
    \node[draw=black, rounded corners, align=center, inner sep=6pt]
      at (-3.8,0.75)
      {\(\partial_\mu T^{\mu\nu}=0\)\\
       \(T^{\mu\nu}=T^{\nu\mu}\)\\
       \(T^\mu{}_\mu=0\)};
    \node[font=\small] at (0.0,2.0) {CKV equation};
    \node[draw=black, rounded corners, align=center, inner sep=6pt]
      at (0.0,0.75)
      {\(\partial_\mu\xi_\nu+\partial_\nu\xi_\mu\)\\
       \(\displaystyle =\frac{2}{d}\delta_{\mu\nu}\partial_\rho\xi^\rho\)};
    \node[font=\small] at (3.9,2.0) {conserved current};
    \node[draw=black, rounded corners, align=center, inner sep=6pt]
      at (3.9,0.75)
      {\(j^\mu_\xi=T^{\mu\nu}\xi_\nu\)\\
       \(\partial_\mu j^\mu_\xi=0\)};
    \draw[-{Latex[length=2mm]}, thick] (-2.15,0.75)--(-1.25,0.75);
    \draw[-{Latex[length=2mm]}, thick] (1.35,0.75)--(2.35,0.75);
    \node[blue!70!black] at (0,-0.9)
      {\(\xi=a,\ \omega x,\ cx,\ b_\rho(2x^\rho x^\mu-x^2\delta^{\rho\mu})\)};
  \end{tikzpicture}
  }
  \caption{A conserved symmetric traceless stress tensor turns every conformal
  Killing vector field into a conserved current.}
  \label{fig:volume-ii-conformal-killing-currents}
\end{figure}

\section{Virial Currents and Improvement}

The fixed-point trace identity determines only the beta-function contribution
to \(T^\mu{}_\mu\).  The stress tensor may still have a total-derivative trace
\[
  T^\mu{}_\mu=\partial_\mu V^\mu
\]
for a local vector operator \(V^\mu\), called a virial current.  In this case
the dilatation current
\[
  D^\mu=x_\nu T^{\mu\nu}-V^\mu
\]
is conserved:
\[
  \partial_\mu D^\mu
  =
  T^\mu{}_\mu-\partial_\mu V^\mu
  =
  0.
\]
This establishes scale invariance generated by \(D^\mu\), but it does not by
itself give conserved special conformal currents of the form
\(T^{\mu\nu}\xi_\nu\).

If the virial current can be removed by an improvement of the stress tensor,
then the same fixed point admits a traceless stress tensor and the conformal
currents above.  For example, if
\[
  V_\mu=(d-1)\partial_\mu L
\]
for a local scalar operator \(L\), then the improvement
\[
  T'_{\mu\nu}
  =
  T_{\mu\nu}
  +
  \bigl(\partial_\mu\partial_\nu-\delta_{\mu\nu}\partial^2\bigr)L
\]
changes the trace by
\[
  T'^\mu{}_\mu
  =
  T^\mu{}_\mu
  -
  (d-1)\partial^2L,
\]
and removes this virial trace.

\begin{remark}[Scale and conformal fixed points]
The implication from scale invariance to conformal invariance depends on the
existence and removability of \(V^\mu\), together with assumptions such as
unitarity, locality, discreteness of scaling dimensions, and the dimension of
spacetime.  The monograph will use the precise stress-tensor criterion above:
conformal symmetry is present when the flat-space theory has a conserved
symmetric stress tensor whose trace can be set to zero by allowed
improvements.
\end{remark}

\begin{figure}[H]
  \centering
  \resizebox{0.9\textwidth}{!}{%
  \begin{tikzpicture}[x=1cm,y=1cm,every node/.style={font=\scriptsize}]
    \node[draw=black, rounded corners, align=center, inner sep=6pt]
      at (-3.2,0.75)
      {\(T^\mu{}_\mu=\partial_\mu V^\mu\)\\
       \(D^\mu=x_\nu T^{\mu\nu}-V^\mu\)};
    \node[draw=black, rounded corners, align=center, inner sep=6pt]
      at (0.75,0.75)
      {\(V_\mu=(d-1)\partial_\mu L\)\\
       improvement exists};
    \node[draw=black, rounded corners, align=center, inner sep=6pt]
      at (4.55,0.75)
      {\(T'^\mu{}_\mu=0\)\\
       \(j^\mu_\xi=T'^{\mu\nu}\xi_\nu\)};
    \draw[-{Latex[length=2mm]}, thick] (-1.55,0.75)--(-0.65,0.75);
    \draw[-{Latex[length=2mm]}, thick] (2.25,0.75)--(3.2,0.75);
    \node[blue!70!black, align=center] at (0.75,-0.65)
      {without such data, scale currents\\do not determine special conformal currents};
  \end{tikzpicture}
  }
  \caption{A virial current gives a conserved dilatation current.  Conformal
  currents follow when the virial trace is removable by an allowed improvement
  of the stress tensor.}
  \label{fig:volume-ii-virial-improvement}
\end{figure}

The next construction searches for nonzero solutions of
\(\beta_J(\lambda)=0\).  The scalar quartic theory in \(4-\epsilon\) dimensions
provides the first perturbative example, after the mass deformation is tuned so
that the infrared theory has no mass gap.
